% P16: source-distinct IBE, contrastive explanation and defeasible appraisal.
\chapter{Inference to the Best Explanation}
\label{sec:p16-inference-to-best-explanation}

An explanation can be a reason to investigate a hypothesis without being
a reason to believe it established. Inference to the best explanation
(IBE) asks how explanatory considerations can contribute to the second
judgment, not merely how a hypothesis can be made compatible with data.
Lipton's 2000 account distinguishes describing our inferential practices
from justifying their reliability; the same distinction organizes this
treatment \cite[author copy pp.~1--4]{Lipton2000IBECompanion}.

The primary presentations cited here are kept distinct: the Companion
article, the two 1993 papers on contrast and underconsideration, and the
1998 discussion of prediction. Specific raven and Bayesian comparisons
are separately attributed to the 2004 second edition. Neither the earlier
papers nor the formal illustrations below are that complete book under
another name. Peirce's historical inquiry stages in
\cref{sec:p16-peirce-abduction-and-inquiry}, the finite logical explanation
relation in \cref{def:p16-abduction-models}, and the 1948 explanation
criterion in \cref{def:p16-dn-potential-explanation} remain separate accounts.

\section{What is inferred, from which alternatives?}
\label{sec:p16-ibe-inferential-task}

\begin{hbdefinition}{def:p16-ibe-task}{An explanatory appraisal problem}
The appraisal specifies observations, a revisable background, the question
to be answered, and a set of considered hypotheses. A \emph{potential
explanation} states how a hypothesis, if correct under the specified
auxiliaries, would account for the observations. Calling it potential does
not assert its truth. An \emph{actual explanation}, in the truth-directed
reading used here, additionally gets the relevant account right.

Competing hypotheses offer incompatible answers to the same question.
Compatible explanations can instead concern different mechanisms, levels
or contributing factors. The appraisal's possible outputs include accepting
a hypothesis provisionally, increasing or decreasing its support,
retaining several compatible explanations, and suspending judgment.
\end{hbdefinition}

The distinction prevents a circular instruction to choose the best already
true explanation. It also prevents forcing a winner whenever candidates
exist. An especially informative false story is still false, and the best
member of a poor collection may not merit acceptance. Lipton accordingly
allows both agnosticism and multiple compatible explanations
\cite[author copy pp.~5--7]{Lipton2000IBECompanion}.

Consider a failed data transfer. A damaged cable and a misconfigured
receiver may be genuine competitors for a specified failure. A weak cable
and poor error recovery can also be compatible contributors. A description
of packet loss and an account of the physical damage need not compete at
all. Comparing them in one undifferentiated contest can remove information
instead of resolving the question. These are stipulated reasoning cases,
not diagnoses of a particular installation.

Candidate generation is consequential. The observation does not deliver
an exhaustive list of possible mechanisms. Background knowledge can suggest
plausible causes, constrain their connection with the data, and exclude
incoherent proposals. It can also conceal a possibility the investigator
has not imagined. Selecting among a fixed list and defending that list
are therefore different achievements. The finite enumeration of
\cref{def:p16-abduction-input} solves the first task for its declared
atomic hypotheses; it does not solve this discovery problem.

\section{Loveliness, likeliness and the guiding claim}
\label{sec:p16-ibe-virtues}

Lipton distinguishes \emph{likeliness}, evidential probability, from
\emph{loveliness}, the understanding an account would provide if correct.
Simply recommending the likeliest hypothesis leaves unexplained how its
likeliness was judged. The explanatory proposal is substantive only if
explanatory features help guide that judgment. His three challenges are
identifying those features, establishing their match with inferential
judgments, and showing that the features actually guide those judgments
\cite[author copy pp.~5--9]{Lipton2000IBECompanion}.

A useful appraisal makes the claimed virtues inspectable. Scope asks which
phenomena are covered and which are left unexplained. Precision asks which
specific outcomes follow under the account. Mechanism asks which causal
connections are supplied rather than merely named. Unification asks whether
apparently separate cases are organized by the same warranted connection.
Simplicity must specify the representation and what complexity is being
counted. None of these questions alone provides a universal numerical
ranking or a probability of truth.

\begin{hbexample}{def:p16-ibe-virtue-comparison}{A mechanism and a relabeling}
Suppose a device shuts down under heavy load. Saying that it has a
heavy-load shutdown tendency reproduces the observed regularity but gives
little information about why it obtains. A proposed thermal cutoff instead
specifies heat production, a sensor threshold, and the action which stops
the device. It offers additional discriminating consequences: changing
cooling should shift the load at which shutdown occurs; the sensor state
should change before the power is cut.

These are conditional predictions under the proposed mechanism, not
observations already obtained. The tendency description and mechanism can
both be correct; their difference in understanding does not make them
incompatible. A rival faulty load sensor can explain the original shutdown
without the predicted temperature dependence. Measuring temperature and
sensor states can therefore compare rivals, while simply preferring a
more elaborate story cannot.
\end{hbexample}

The guiding claim is not established by listing attractive virtues.
Correlation between explanatory and evidential judgments could have another
explanation. Lipton's descriptive replies appeal to the explanatory
account's coverage of scientific inferences and to its account of why
the two judgments match. These are arguments about the account, not
controlled psychological demonstrations or a solution to wholesale
inductive skepticism \cite[author copy pp.~8--11]{Lipton2000IBECompanion}.

\section{The causal question is contrastive}
\label{sec:p16-ibe-contrast}

In \emph{Making a Difference}, Lipton asks which causes explain a chosen
contrast. The fact occurs and the foil does not, but the two descriptions
need not be logically incompatible. Two people could both receive a prize;
asking why one did rather than the other does not require that the first
person's success prevented the second's. Changing the foil can change
which causal difference is relevant \cite[secs.~1--2, author copy
pp.~1--8]{Lipton1993MakingDifference}.

The exclusion account---explain the fact and then deduce the foil's
absence---misses compatible contrasts. Lewis's proposed cause that would
not have caused the foil is also insufficient: a cause located in the
first instance may not cause the second instance, even when the second
has a corresponding cause of the same relevant type. Lipton instead
investigates a cause of the fact with no corresponding token in the foil
instance. This is not mere difference counting: correspondence, causal
type, alternative differences and nonactual foils all need an account
\cite[sec.~2]{Lipton1993MakingDifference}.

\begin{hbdefinition}{def:p16-ibe-printer}{A stipulated contrast model}
Let the binary variables \(P,R,T,J\) specify power, paper, usable toner
and a jam. Define awake state \(A=P\) and successful output
\[
                    O=P\land R\land T\land\neg J.
\]
These are total deterministic assignments. The value one denotes presence
of the named condition. In context \(u\), put \(P=R=T=J=1\).
In context \(v\), put \(P=R=T=1,J=0\).
\end{hbdefinition}

\begin{hbtheorem}{thm:p16-ibe-printer}{proved=annex:p16.ibe-printer}{A specified difference and a bad remedy}
The contexts give \((A,O)=(1,0)\) and \((1,1)\), respectively.
Replacing \(J\) by zero in context \(u\) produces successful output
while retaining the awake state. Replacing \(P\) by zero removes output
and also the awake state. Power is present on both sides of the chosen
output contrast; citing it does not distinguish these two contexts.
\end{hbtheorem}

The calculation supplies the entire local mechanism and the unchanged
background. It does not establish from two observations alone that a jam
was the actual cause. A mistaken paper sensor could produce the same
recorded symptom. Nor does eliminating a necessary condition for operation
count as explaining failure relative to normal operation. The intended
foil controls the diagnostic task.

For nonactual foils, Lipton investigates how a contrast focuses the
alternative outcome and fixes shared portions of the causal history.
An actual comparison instance can serve as a surrogate for a question
about the original instance. He distinguishes this backward-counterfactual
question from an intervention which deliberately changes an earlier
equation. His objections to holding the entire past fixed at the cost
of a last-minute law violation, and his proposed focusing/fixing response,
are philosophical arguments, not a general structural counterfactual
identification theorem \cite[secs.~3--5]{Lipton1993MakingDifference}.

The formal semantics of \cref{def:p16-causal-intervention} can specify a
particular manipulation. They do not uniquely specify Lipton's relevant
foil, correspondence relation or resolution of a backward counterfactual.
The causal account also does not cover every explanation: logical and
mathematical explanations need not cite a physical cause.

\section{Contrastive inference and rival stories}
\label{sec:p16-ibe-comparative-case}

Lipton's Companion account uses Semmelweis's investigation to distinguish
an account of fever from an account of its distribution between maternity
wards and before/after disinfection. Conditions common to both wards do
not explain the ward contrast without additional premises; the proposed
infection mechanism addresses both comparisons. This is the source's
historical illustration, not a newly reconstructed dataset, experiment
or clinical recommendation \cite[author copy p.~8]{Lipton2000IBECompanion}.

\begin{hbexample}{def:p16-ibe-transfer-case}{A two-stage failure investigation}
Suppose two otherwise comparable transfer stations show different failure
rates. The initial candidates are network congestion, a common software
release and a station-specific damaged connector. Both stations use the
release and share the same upstream network. The connector proposal
addresses the station contrast; the release and network proposals need
an additional station-specific interaction to do so.

Now replace the connector, retain the release and network, and observe
that only the formerly failing station improves. The connector proposal
also addresses this temporal contrast. Its advantage is conditional on
the comparison being genuine: the replacement must not also change an
unrecorded driver setting, workload or reporting policy. If those changes
occurred together, the temporal contrast does not distinguish them.

A previously unconsidered connector-firmware interaction is a serious
defeater if failures recur only with one firmware version. It does not
erase the evidence for the connector's involvement, but it defeats the
stronger account that connector condition alone suffices. The revised
question compares faulty and sound connectors within each firmware
condition. Additional observations can refine rather than simply replace
the explanation.
\end{hbexample}

The substantive work is locating mechanisms which address the actual
contrasts, assessing alternatives, and defending the shared history and
measurement assumptions. The following mathematical bridges make certain
additional premises precise; they do not convert explanatory appeal into
deductive necessity.

\section{The raven problem: logic is not a sampling policy}
\label{sec:p16-ibe-ravens}

\begin{hbtheorem}{thm:p16-ibe-raven-logic}{proved=annex:p16.ibe-raven-logic}{Equivalence and the limit of one favorable instance}
For arbitrary predicates \(R,B\) on a nonempty domain,
\[
 \forall x(Rx\to Bx)\quad\text{and}\quad
 \forall x(\neg Bx\to\neg Rx)
\]
are classically equivalent. An object with \(\neg R\land\neg B\)
satisfies both instances. On a domain with at least two elements, that
observation alone does not entail either universal: another object may
have \(R\land\neg B\).
\end{hbtheorem}

Lipton's second-edition response distinguishes suitable from irrelevant
contrapositive instances. Contrastive support needs comparable histories
and no known preempting explanation of the foil's absence. Agreement
instead varies suspect background conditions while retaining a credible
common etiology. A shoe is not a suitable fever comparison. His discussion
does not claim that contrastive inference itself establishes the raven
generalization \cite[pp.~95--101]{Lipton2004BestExplanation}.

\begin{hbdefinition}{def:p16-ibe-raven-sampling}{Two explicitly stipulated populations}
There are 100 labeled objects, exactly two of them ravens and 98 white
nonravens. Hypothesis \(H_0\) makes both ravens black; \(H_1\) makes one
black and one nonblack. Give these competing population hypotheses equal
prior probability. In protocol U choose uniformly from all 100 objects.
In protocol V choose uniformly from the two ravens. In this probability
calculation, \(B\) denotes the event that the selected object is black,
rather than the predicate symbol in the preceding logical statement.
The selection protocol
is known and independent of the hypothesis conditional on the specified
population. This is an artificial finite model, not a claim about birds.
\end{hbdefinition}

\begin{hbtheorem}{thm:p16-ibe-raven-sampling}{proved=annex:p16.ibe-raven-sampling}{Relevance depends on the declared observation model}
Under U, observing a white nonraven leaves \(P(H_0)=1/2\).
Under V, observing a black raven gives \(P(H_0\mid B)=2/3\).
The observation in the second case confirms but does not entail \(H_0\).
\end{hbtheorem}

Equivalence of the universal descriptions does not select U or V, their
population counts or their priors. A different observation model could
make a nonraven informative through an explicitly defended dependence.
The calculation is not a theorem that every white object is irrelevant,
nor a numerical reconstruction of Lipton's causal argument.

\section{What a Bayesian bridge adds}
\label{sec:p16-ibe-bayes}

The second edition proposes compatibility between explanatory thinking
and Bayesian updating, not an extra explanatory bonus attached after
conditionalization. Explanatory considerations can concern candidate
construction or the assessment of the probabilistic inputs
\cite[pp.~107--108]{Lipton2004BestExplanation}. This does not identify
IBE with posterior maximization: the probability calculus does not itself
choose the hypotheses, their explanatory role, the prior or the likelihood.

Use the finite conditional-probability prerequisites in
\cref{sec:p15-finite-conditioning,sec:p15-finite-independence}.
The general updating/prediction treatment is
\cref{sec:p15-bayesian-updating-prediction}. No numerical parameter in
the following definition is licensed merely by calling an account lovely.

\begin{hbdefinition}{def:p16-ibe-bayes-input}{An explicitly additional finite model}
On one probability space let \(B_0\) be background evidence with positive
probability. Events \(H_1,\ldots,H_m\) partition the space, \(m\ge2\),
and have positive conditional prior masses
\(\pi_i=P(H_i\mid B_0)\). Let \(E\) be new evidence and put
\(\ell_i=P(E\mid H_i,B_0)\), with
\(Z=\sum_i\pi_i\ell_i>0\). A nonempty considered subset
\(I\subseteq\{1,\ldots,m\}\) has event
\(C_I=\bigcup_{i\in I}H_i\) and positive posterior mass \(r\).
These exhaustive model events are distinguished from the investigator's
possibly incomplete considered subset.
\end{hbdefinition}

\begin{hbtheorem}{thm:p16-ibe-bayes-scope}{proved=annex:p16.ibe-bayes-scope}{Posterior odds and the bad-lot multiplier}
Write \(a_i=P(H_i\mid E,B_0)\). Then
\[
 a_i=\frac{\pi_i\ell_i}{Z},\qquad
 \frac{a_i}{a_j}=\frac{\pi_i}{\pi_j}\frac{\ell_i}{\ell_j}
 \quad(\ell_j>0).
\]
For \(i\in I\), write \(q_i=P(H_i\mid E,B_0,C_I)\). Then
\[
                 a_i=rq_i,\qquad r=\sum_{i\in I}a_i.
\]
Thus being the highest-ranked considered hypothesis, even with a large
\(q_i\), supplies no lower bound on its absolute posterior independent
of \(r\). For posterior masses \((.18,.02,.80)\) and
\(I=\{1,2\}\), \(q_1=.90\) but \(a_1=.18\).
\end{hbtheorem}

The example can be realized with evidence certain and those three prior
masses; its purpose is the ranking/coverage distinction, not confirmation
by a particular experiment. An investigator who conditions away the third
possibility without justification has changed the question. A numerical
acceptance threshold would add another normative premise; it is not supplied
by the displayed maximum.

\begin{hbtheorem}{thm:p16-ibe-bonus}{proved=annex:p16.ibe-bonus}{A nonconstant bonus changes a specified posterior}
Suppose all posterior masses \(a_i\) are positive, and choose positive
weights \(b_i\). The alternative vector
\[
                 a_i^*=\frac{b_i a_i}{\sum_jb_j a_j}
\]
equals the specified Bayesian posterior at every index exactly when the
\(b_i\) are constant. This algebra does not prohibit revising a prior,
likelihood or evidence model for an independently justified reason. It
does show that a nonconstant bonus is not the unchanged conditionalization.
\end{hbtheorem}

\section{Underconsideration, background and actual replies}
\label{sec:p16-ibe-bad-lot}

Lipton's 1993 argument addresses a specific intermediate skepticism:
known reliable comparison of considered theories is conceded, while any
warrant for having generated a true candidate is denied. This is stronger
than the harmless observation that candidates are fallible. His replies
use contradictory alternatives and the background theories on which
actual local evaluation depends; they do not claim to defeat unrestricted
Humean skepticism \cite[secs.~I--VII]{Lipton1993BestGoodEnough}.

The contradictory-pair reply distinguishes a theory and its negation from
two merely incompatible detailed theories. If a reliable ranking genuinely
covers a proposition and its negation, the pair exhausts the possibilities.
The skeptic can restrict the ranking concession, or contend that only the
uninformative negation will win, but those replies change the conceded
inductive powers or the claim that generation never reaches interesting
truths. The probability version below shows exactly what its premises buy.

\begin{hbtheorem}{thm:p16-ibe-contradictories}{proved=annex:p16.ibe-contradictories}{A conditional ranking-to-threshold implication}
Under any probability law, \(P(H)>P(\neg H)\) implies \(P(H)>1/2\).
If \(K\subseteq\neg H\) and \(P(K)>P(H)\), then
\(P(\neg H)>P(H)\). These implications require the stated probability
comparisons; an unexplained aesthetic ranking does not supply them.
Equality of a contradictory pair gives no strict threshold conclusion.
\end{hbtheorem}

The background reply is more substantive. Instruments, data descriptions,
candidate plausibility and evidential bearing are assessed through already
accepted theories. Lipton argues that granting sustained reliable local
ranking while allowing this background to be arbitrarily wrong is unstable;
iteration makes today's favored theories tomorrow's evaluative background.
He extends the objection to certain arguments from permanent inductive ties,
and challenges a global ranking which removes the actual local background
from view \cite[secs.~V--VII]{Lipton1993BestGoodEnough}.

This argument depends on a description of real evaluative dependence and a
truth connection for that background. It is not a logical theorem that an
arbitrary ranking device needs every background sentence true. A device
could ignore a false input, and an approximate physical theory could support
some accurate comparisons while failing elsewhere. The required claim is
about relevant background and the warranted range of its use. Lipton's
own conclusion is fallibilist: he does not establish perfect ranking,
defeat every intermediate skeptic, or prove that truth rather than empirical
adequacy is required for every reliable assessment.

The objection is also not uniquely directed at explanatory reasoning.
An incomplete candidate space can compromise a nonexplanatory comparison.
Conversely, redescribing IBE's output as empirical adequacy does not remove
underconsideration: adequacy to everything observable is stronger than fit
to the observations already collected. This preserves the source's distinction
between realist commitments and the form of an explanatory appraisal.

\section{Prediction, accommodation and selection history}
\label{sec:p16-ibe-prediction}

Lipton's 1998 discussion separates evidence known before a theory is
constructed from evidence subsequently checked. Calling an accommodation
ad hoc is either merely saying it was purpose-built, or assuming the poor
support one was meant to establish. Likewise, the theorist's risk of a
publicly failed prediction is not automatically a difference in how a fixed
theory bears on a fixed observation. His positive proposals concern the
choice of discriminating tests and the opportunity to alter a theory to fit
known results \cite[author copy pp.~10--15]{Lipton1998CambridgeContributions}.

\begin{hbdefinition}{def:p16-ibe-selection-input}{A finite selection illustration}
Let \(X\) be a uniformly distributed binary string of length \(k\ge1\).
Before seeing it, an investigator fixes a string \(s\). A second
procedure observes \(X\) and then declares the chosen string to be
\(S(X)=X\). Both report whether the chosen string equals the observation.
The two procedures, not just the final matching string, are specified.
\end{hbdefinition}

\begin{hbtheorem}{thm:p16-ibe-selection-fit}{proved=annex:p16.ibe-selection-fit}{A selected fit is a different event}
The fixed prediction succeeds with probability \(2^{-k}\), whereas
the post-observation selected match succeeds with probability one.
For \(k=8\), the two probabilities are \(1/256\) and one.
This does not assign a probability of truth to a scientific theory, or
prove that all accommodation is defective.
\end{hbtheorem}

The event in the second procedure is not \(X=s\) for a beforehand fixed
\(s\). The selected answer can fit perfectly while revealing nothing about
whether the procedure will fit new data. But a mechanism genuinely
discovered from old data can also make independently constrained and
successful new predictions. Merely knowing the dates of theory and
observation does not determine the opportunity for selection or the
strength of the constraints. Prediction may fail to provide an advantage
when tests are chosen after seeing their results or unfavorable predictions
are suppressed; accommodation may be informative when its constraints
were independently fixed.

The proposed fudging explanation is therefore a reason to inspect the
construction process, not permission to discount every accommodated
observation or award every novel prediction a fixed multiplier. Background
assumptions about which adjustments were available, which were made and
what they would predict must themselves be defended. The finite example
exposes one such dependency without presenting it as Lipton's general theorem.

\section{Defeaters and reference use}
\label{sec:p16-ibe-reference-use}

A rebutting defeater supplies grounds against the proposed conclusion.
An undercutting defeater attacks the link from evidence to conclusion:
a faulty instrument, an unsuitable comparison, copied testimony or a
selection procedure can weaken the inference without showing the
hypothesis false. These operational categories organize the illustrative
appraisal; they are not asserted as Lipton's own formal calculus.

\begin{hbexample}{def:p16-ibe-failed-investigation}{When the lovely story survives by moving its premises}
The thermal-cutoff investigator first predicts rising temperature before
shutdown. The temperature record stays normal. She then declares the
sensor inaccurate, later proposes an unobserved local hotspot, and finally
changes the threshold after every new result. Each modification may be
possible, but possibility is not independent support. If the auxiliaries
are unconstrained, the account no longer discriminates the rival load-sensor
mechanism or says which observation would weaken it.

The appropriate response is to record the failed consequence and the
changed auxiliaries, test the sensor and hotspot proposals independently,
and retain suspension of judgment as an option. Rejecting a faulty
temperature sensor would not by itself establish the thermal mechanism;
confirming a genuine hotspot would not make every subsequent adjustment
legitimate. The explanatory account and its measurement assumptions remain
separate objects of appraisal.
\end{hbexample}

\begin{hbmethod}{def:p16-ibe-consultation}{Consulting an explanatory inference}
\begin{enumerate}
\item Fix the evidence, background, explanatory question and any foil.
      Distinguish observations from assumptions which generated their description.
\item Record considered rivals and compatible contributors. Ask what the
      search could have omitted; do not equate the considered set with exhaustion.
\item State what each candidate would explain and which consequences depend
      on which auxiliaries. Make the claimed virtue and comparison inspectable.
\item Seek observations which distinguish candidates at the chosen scope.
      Establish shared-history, measurement and selection premises explicitly.
\item Separate explanatory preference, increased support, provisional
      acceptance and truth. Any numerical or decision threshold needs its own warrant.
\item Appraise rebutting and undercutting evidence, retain failures, and revise
      the model or suspend judgment when the comparison no longer discriminates.
\end{enumerate}
\end{hbmethod}

Successful use is a conditional, revisable appraisal with actual alternatives
and warranted connections. Failure includes a best-of-a-bad-list inference,
an irrelevant foil, a relabeled symptom, a posterior computed after excluding
an unlicensed alternative, or a hypothesis protected by endlessly moving
auxiliaries. None is repaired by treating explanatory preference as
deductive entailment. The logical, structural and probability results
linked above make specified neighboring tasks precise; they do not
silently settle the philosophical warrant of IBE.

\index{inference to the best explanation}
\index{explanation!contrastive}
\index{explanation!loveliness and likeliness}
\index{underconsideration}
\index{raven paradox}
\index{prediction!accommodation and selection}
